\section{Introduction and scientific scope}

\subsection{Motivation}

Previous waveform-based machine-learning studies have shown that learned models can improve timing in time-of-flight positron emission tomography (TOF-PET), including convolutional neural network (CNN)-based TOF estimation and explicit timing-correction approaches~\cite{berg2018,onishi2022,feng2024,naunheim2025}. The present study focuses on a narrower question: whether the useful event-by-event correction requires a paired convolutional representation, or whether it can be captured by a detector-level function shared between the two channels.

The final analysis therefore does not present a broad model benchmark. It compares only the two models retained in the waveform pipeline: the proposed antisymmetric multilayer perceptron (MLP) and the paired Onishi CNN reference. Both receive the same synchronized detector-pair waveforms and predict a correction to the calibrated leading-edge-discrimination (LED) timing residual.

\subsection{Project scope}

The proposed model applies one shared detector scorer \(g_\theta\) independently to the two detector waveforms and forms the pair correction as
\begin{equation}
    y_\theta(s_1,s_2)=g_\theta(s_1)-g_\theta(s_2).
\end{equation}
This representation enforces exact detector-exchange antisymmetry while avoiding a convolutional prior inside the detector scorer.

The reference model is the paired CNN based on Onishi et al.~\cite{onishi2022}. It processes the two detector waveforms jointly, with the first two-dimensional convolution spanning the detector axis, and predicts one pair-level correction.

The comparison is repeated for two fixed waveform regions defined in the shared analysis profile:
\begin{itemize}
    \item the \textbf{Onishi window}, from \SI{-1.5}{ns} to \SI{2.0}{ns};
    \item the \textbf{wide window}, from \SI{-2.0}{ns} to \SI{30.0}{ns}.
\end{itemize}
This separates the model-architecture question from the amount of temporal context provided to the model.

The same framework can be applied to the first-stage energy waveform and to the dedicated timing waveform. The two waveform families answer different practical questions: the energy signal may leave more correctable time walk and pulse-shape dependence, whereas the dedicated timing signal can provide a substantially better conventional timing baseline.

\subsection{Research questions}

The report is organized around three primary questions:

\begin{enumerate}[label=\textbf{Q\arabic*.}]
    \item \textbf{Model representation:} how does the detector-shared antisymmetric MLP compare with the paired Onishi CNN on the same blind-test events?

    \item \textbf{Temporal context:} does the relative behavior of the two models change when the input is restricted to the Onishi window or extended to the wider waveform region?

    \item \textbf{Signal and hardware dependence:} are the same conclusions reproduced for energy- and timing-channel waveforms and, where available, across the UC and FBK front-end boards?
\end{enumerate}

LED-threshold scans, model-output correlation, and gradient-based explainability are treated as supporting analyses. They help interpret the two-model comparison but do not introduce additional competing models.

\reportfigure
    {figures/paper_scope.pdf}
    {Scientific scope of the final study: antisymmetric MLP versus paired Onishi CNN, evaluated with two waveform windows and repeated across the available waveform families and front-end datasets.}
    {fig:paper-scope}
